\documentclass[11pt,a4paper]{article}
\usepackage[margin=16mm,headheight=15pt]{geometry}
\usepackage{fontspec}
\usepackage{polyglossia}
\setdefaultlanguage{russian}
\setotherlanguage{english}
\setmainfont{Arial}
\setsansfont{Arial}
\setmonofont{Menlo}
\newfontfamily\cyrillicfont{Arial}
\newfontfamily\cyrillicfonttt{Menlo}
\usepackage{unicode-math}
\setmathfont{STIX Two Math}
\usepackage{xcolor,enumitem,amsmath}
\usepackage[most]{tcolorbox}
\usepackage{fancyhdr}
\usepackage[colorlinks=true,linkcolor=blue,urlcolor=blue]{hyperref}
\definecolor{navy}{HTML}{102044}
\definecolor{accent}{HTML}{1168E8}
\definecolor{pale}{HTML}{F3F8FF}
\definecolor{ink}{HTML}{33475F}
\definecolor{sand}{HTML}{FFF6D8}
\definecolor{goldline}{HTML}{C49200}
\setlength{\parindent}{0pt}
\setlength{\parskip}{0.35em}
\emergencystretch=3em
\pagestyle{fancy}
\fancyhf{}
\fancyhead[L]{\small\color{navy} Пара 2 · что говорить}
\fancyhead[R]{\small\color{accent} 34 слайда}
\fancyfoot[C]{\small\thepage}
\renewcommand{\headrulewidth}{0.4pt}
\newtcolorbox{saybox}{colback=pale,colframe=accent,boxrule=0pt,leftrule=3.5pt,arc=1.5mm,breakable,left=3.5mm,right=3.5mm,top=2mm,bottom=2mm,before upper={\fontsize{12.5}{16.5}\selectfont}}
\newcommand{\slidehead}[2]{\par\vspace{4mm}{\color{navy}\large\bfseries Слайд #1. #2}\par\vspace{1mm}}
\newcommand{\screen}[1]{\par{\color{ink}\small #1\par}\vspace{1mm}}
\newcommand{\say}[1]{\begin{saybox}#1\end{saybox}}
\newtcolorbox{signbox}{colback=sand,colframe=goldline,boxrule=0pt,leftrule=3.5pt,arc=1.5mm,breakable,left=3.5mm,right=3.5mm,top=2mm,bottom=2mm,before upper={\fontsize{12}{15.5}\selectfont}}
\newenvironment{signs}{\begin{signbox}{\bfseries Что значит каждый знак на этом слайде.}\par\begin{itemize}[leftmargin=1.15em,itemsep=2pt,topsep=3pt,parsep=0pt]}{\end{itemize}\end{signbox}}
\begin{document}

{\color{navy}\bfseries\fontsize{24}{30}\selectfont Что говорить и что значит каждый знак\par}
\vspace{2mm}
Презентация: \url{http://localhost:8000/materials/Пара2/?v=spot24}.\\
Обновите страницу через Cmd+Shift+R. Слайдов 34.

На каждом слайде три слоя.

Серая строка~--- что на экране. Студентам не читай.

Жёлтая рамка «Что значит каждый знак»~--- только там, где на слайде формула. Слева знак, как на экране. Справа слова, что он значит. Это разбор формулы.

Голубая рамка~--- слова, которые говоришь студентам. Жёлтую полосу на самом слайде говори медленнее.

Файл, который надо открыть: \texttt{output/pdf/МЕТОДИЧКА\_ПРЕПОДАВАТЕЛЯ.pdf}. Если открыт старый~--- закрой и открой этот заново.

Систему на доске не решай. Готовые числа называй, когда дойдёшь до метрик: старт прямой 0,67, наклон 0,5, доля разброса 0,75.

Если к 60-й минуте ты ещё на нормальных уравнениях, на слайде 15 прочитай только жёлтую полосу и переходи к метрикам. Вслух обязательно четыре: средняя квадратичная ошибка, корень из неё можно назвать, средняя абсолютная ошибка, доля разброса и процентная ошибка. Остальные метрики назови и перелистни.

\slidehead{1}{Титул}
\screen{Заголовок «Линейная регрессия». Сегодня ищем веса и смотрим, какое число можно честно назвать качеством.}
\say{Сегодня линейная регрессия. Прогноз~--- это сдвиг плюс признаки, умноженные на веса. Веса будем уменьшать ошибку шаг за шагом, и есть формула того места, куда эти шаги приходят. Потом разберём, какое число после этого можно назвать качеством, а какое получилось просто потому, что мы посчитали его на тех же строках, где искали веса.}

\slidehead{2}{Вопрос пары}
\screen{Два вопроса на всю пару. Откуда веса. И какое число можно назвать качеством.}
\say{Два вопроса на всю пару.

Первый. Откуда берутся веса. Вес~--- это число при признаке: на сколько рублей меняется прогноз, если площадь выросла на один метр. Ищем их так: на каждом шаге ошибка становится меньше, пока не сядет на ноль. Формула записывает то место, куда шаги пришли. Картинку шагов увидите на слайде с нормальными уравнениями.

Второй. Какое число после этого можно назвать качеством. Красивый счёт на тех же строках, где искали веса,~--- это ещё арифметика. Качеством называем число, посчитанное на других строках.}

\slidehead{3}{Гальтон}
\screen{Три строки про рост родителей и сыновей. Внизу английская фраза: возврат к середине.}
\say{Слово «регрессия» пришло из опыта про рост. В 1885--1886 годах Фрэнсис Гальтон сравнивал рост родителей и рост взрослых сыновей.

Три строки. Высокие родители: сыновья тоже высокие, но в среднем ниже отцов. Низкие родители: сыновья тоже низкие, но в среднем выше отцов. Средние родители: сыновья остаются около середины.

Гальтон назвал это возвратом к середине. На слайде по-английски: regression towards mediocrity. Отсюда имя метода.

Если спросят про квадраты. Метод наименьших квадратов старше Гальтона: его опубликовал Лежандр в 1805 году. Пирсон позже записал корреляцию. Квадраты он не придумал.}

\slidehead{4}{Картинка роста}
\screen{Мелкие точки~--- семьи. Крупные~--- средний сын в группе. Прямая положе диагонали, наклон около двух третей.}
\begin{signs}
\item \(E\) — среднее. Не ошибка и не сумма.
\item \(y\) — рост сына, ответ.
\item \(x\) — рост родителей, признак.
\item \(E[y\mid x]\) — средний рост сына при таком росте родителей. Вертикальная черта читается «при условии».
\item \(\mu\) — середина популяции, средний рост.
\item \(\beta\) — доля отклонения родителей, которая переходит сыну. При близких разбросах это корреляция. Меньше 1, если связь не идеальна.
\item \(x-\mu\) — насколько родители выше или ниже середины.
\item Плюс перед \(\beta\) — к середине прибавляем эту долю, не весь рост родителей.
\end{signs}
\say{Мелкая точка~--- одна семья. Крупная~--- средний рост сыновей у родителей похожего роста. Крупные точки легли на прямую, и она положе диагонали.

Формула справа, слева направо. Буква \(E\)~--- среднее. В скобках \(y\) при условии \(x\): средний рост сына при таком росте родителей. Равно. \(\mu\)~--- середина популяции. Плюс \(\beta\), умноженная на скобку. В скобке \(x-\mu\): насколько родители выше или ниже середины. \(\beta\)~--- доля этого отклонения, которая переходит сыну. При близких разбросах это корреляция, и она меньше 1, если связь не идеальна.

У Гальтона \(\beta\) около двух третей. Родители на 6 сантиметров выше середины, сыновья в среднем примерно на 4 сантиметра выше. Диагональ отдала бы сыну все 6.

Крайние родители не передают детям случайный кусок своего роста, поэтому следующая группа ближе к середине. Тот же эффект без биологии: отобрали худших по первой контрольной, на второй они в среднем ближе к общей середине, даже если курс ничего не изменил.}

\slidehead{5}{Что такое регрессия}
\screen{Ответ~--- число. При тех же признаках ответы всё равно разбросаны. Сегодня в прогноз пишем среднее.}
\begin{signs}
\item \(E[y\mid x]\) — средний ответ при этих признаках. Это сегодняшняя сводка, её и пишем в прогноз.
\item \(x^{\top}\) — признаки одной квартиры, записанные строкой. Единица для сдвига уже внутри этой строки.
\item \(w\) — столбец весов.
\item \(x^{\top}w\) — каждый признак умножили на свой вес и сложили.
\item Знак \(=\) — среднее ответа и есть эта взвешенная сумма.
\end{signs}
\say{Регрессия здесь предсказывает число: цену, рост, спрос. Метку «да или нет» она не ставит. Если в названии есть слово regression, а ответ~--- класс, это другая задача. Её разберём на следующей паре.

При одних и тех же признаках ответы всё равно разные. Регрессия выбирает, какую сводку этого разброса записать в прогноз. Сегодня сводка~--- среднее.

Третья строка слайда~--- формула. Слева \(E[y\mid x]\): средний ответ при этих признаках. Справа \(x^{\top}w\): признаки, умноженные на веса и сложенные. В прогноз пишем именно это среднее. Один признак~--- прямая. Много признаков~--- плоскость.}

\slidehead{6}{Четыре названия}
\screen{Четыре строки. Сегодня первые две. В обеих на выходе одна цена.}
\say{Четыре названия. Сегодня первые две. В обеих на выходе одно число.

Простая. Одна цена по одному признаку. Двадцать плюс два умножить на площадь. Сорок метров: двадцать плюс восемьдесят, это сто.

Множественная. Признаков несколько, цена всё равно одна. Двадцать плюс два на площадь плюс три на этаж. Сорок метров и пятый этаж: сто плюс пятнадцать, сто пятнадцать. Признаков два, ответ один. В этом вся разница.

Полиномиальная. Среди признаков есть квадрат площади, график гнётся. Считается тем же способом. Сегодня не трогаем.

Логистическая. Ответ уже не цена, а «купит» или «не купит». Следующая пара.}

\slidehead{7}{Формула прогноза}
\screen{Прогноз равен сдвигу плюс веса признаков. Под формулой три слова: вес, сдвиг, остаток.}
\begin{signs}
\item \(\hat y\) — прогноз. Шляпка значит «посчитали моделью», это не настоящая цена.
\item \(y\) без шляпки — настоящий ответ.
\item \(w_0\) — сдвиг. Прогноз, когда все признаки нули.
\item \(w_1\) — вес первого признака: на сколько меняется прогноз, если этот признак увеличить на 1, остальные не трогать.
\item \(x_1\) — первый признак, например площадь.
\item \(\cdots\) — остальные признаки пишутся так же: свой вес на свой признак.
\item \(w_d\) и \(x_d\) — вес и признак последнего. Буква \(d\) — сколько признаков.
\item \(w^{\top}x\) — та же сумма одной строкой. Значок \(\top\) ставит веса в строку, чтобы их можно было умножить на столбец признаков \(x\).
\item \(w_j\) под формулой — вес признака с номером \(j\).
\item остаток \(y-\hat y\) — настоящий ответ минус прогноз. Сколько модель не угадала.
\end{signs}
\say{Формула сверху, слева направо.

Шляпка над \(y\)~--- прогноз. Без шляпки \(y\)~--- настоящий ответ.

Равно \(w_0\). Это сдвиг: прогноз, когда все признаки нули. В квартире это была двадцатка.

Плюс \(w_1 x_1\)~--- первый признак на свой вес. Вес~--- на сколько меняется прогноз, если этот признак увеличить на единицу, а остальные не трогать.

Многоточие~--- так же для остальных признаков. Плюс \(w_d x_d\)~--- последний.

Второе равно, \(w^{\top}x\)~--- та же сумма одной строкой. Чтобы сдвиг туда поместился, в \(x\) стоит единица.

Под формулой три подписи. \(w_j\)~--- вес. \(w_0\)~--- сдвиг. Остаток~--- \(y-\hat y\): сколько модель на этой квартире не угадала.}

\slidehead{8}{Матрица}
\screen{Таблица квартир. Первый столбец~--- единицы. Внизу пример: 1, 40 и 5 умножаются на 20, 2 и 3. Получается 115.}
\begin{signs}
\item \(\hat y\) — столбец прогнозов всех квартир сразу.
\item \(X\) — таблица. Одна строка — одна квартира. Один столбец — один признак.
\item первый столбец из единиц — не из данных. Единица умножается на сдвиг. Без неё нулевая площадь стоила бы ноль.
\item точки в матрице — «остальные строки такие же», не пропущенные числа.
\item \(w\) — столбец весов, в том же порядке, что столбцы таблицы.
\item \(Xw\) — каждую строку умножили на веса и сложили. Это и есть прогноз.
\item нижняя строка \(1,\ 40,\ 5\) — одна квартира: единица, 40 метров, 5-й этаж.
\item столбец \(20,\ 2,\ 3\) — веса: сдвиг 20, вес площади 2, вес этажа 3.
\item \(1\cdot 20+40\cdot 2+5\cdot 3=115\) — та же квартира, что на прошлом слайде.
\end{signs}
\say{Две формулы.

Верхняя. \(\hat y=Xw\). \(X\)~--- таблица квартир, \(w\)~--- столбец весов. Одна строка таблицы~--- одна квартира.

Смотри на матрицу \(X\). Первый столбец~--- единицы, они не из данных. Единица умножается на сдвиг. Без неё квартира с нулевой площадью обязана стоить ноль. Дальше столбцы признаков. Точки сверху вниз~--- остальные строки такие же. Строк столько, сколько квартир.

Нижняя строка, слева направо. Строка \(1,\ 40,\ 5\)~--- одна квартира: единица для сдвига, 40 метров, 5-й этаж. Столбец \(20,\ 2,\ 3\)~--- веса. \(1\cdot 20=20\). \(40\cdot 2=80\). \(5\cdot 3=15\). Сумма 115. Это та же квартира, что на прошлом слайде: 20 плюс 80 плюс 15.

В программе столбец единиц добавляется сам. Параметр называется fit intercept, по умолчанию он включён.}

\slidehead{9}{One-hot и one-cold}
\screen{Таблица трёх районов. Слева одна единица в строке. Справа один ноль в строке. Коды 1, 2, 3~--- ни то ни другое.}
\say{Район~--- это название, не число. Если записать Центр как 1, Север как 2, Юг как 3, модель прочитает Юг как втрое больший Центр. Так нельзя.

Слева one-hot. В строке одна единица, она отмечает свой район. Центр: 1 0 0. Север: 0 1 0. Юг: 0 0 1. Покажи строку «Север»: единица посередине.

Справа one-cold. Это единица минус левая таблица. У Севера ноль стоит посередине, а по краям единицы: 1 0 1.

Оба варианта~--- обычные столбцы рядом с площадью. Почему вместе со сдвигом они ломают веса, скажем через один слайд.}

\slidehead{10}{Картинка районов}
\screen{Те же три района клетками. Синяя клетка~--- единица, бледная~--- ноль.}
\say{Та же таблица, только клетками. Синяя~--- единица, бледная~--- ноль. Слева one-hot, справа one-cold. Найди Север: слева синяя середина, справа в середине ноль, по краям единицы. Новой мысли нет, картинка показывает замену нуля и единицы.}

\slidehead{11}{Почему веса не единственные}
\screen{Слева сумма трёх столбцов one-hot равна столбцу единиц. Справа сумма one-cold равна двум таким столбцам. Жёлтая полоса: сдвиг оставляем, один район выбрасываем.}
\begin{signs}
\item \(1_{\text{ц}}\), \(1_{\text{с}}\), \(1_{\text{ю}}\) — три столбца one-hot. В столбце единица стоит только у своего района.
\item \(\mathbf{1}\) — столбец сдвига: единица в каждой строке. Жирная единица, не число «район номер 1».
\item сумма слева \(=\mathbf{1}\) — сложи три столбца районов, получится столбец сдвига. Сдвиг уже спрятан внутри районов.
\item \(c_{\text{ц}}\), \(c_{\text{с}}\), \(c_{\text{ю}}\) — три столбца one-cold. Буква \(c\) — «холодный» код, ноль в своём районе.
\item \(2\mathbf{1}\) — два столбца сдвига. В каждой строке one-cold две единицы, поэтому сумма столбцов вдвое больше столбца единиц.
\end{signs}
\say{Нельзя оставить и сдвиг, и все три района сразу.

Слева формула one-hot: \(1_{\text{ц}}+1_{\text{с}}+1_{\text{ю}}=\mathbf 1\). Три единичных столбца районов в сумме дают жирную единицу. Жирная единица~--- столбец сдвига. В каждой строке ровно одна единица района, поэтому сдвиг уже спрятан внутри районов.

Один и тот же прогноз можно записать разными весами. Хотим Юг 100, Центр 120, Север 90. Первый способ: Юг выбрасываем, сдвиг 100, Центр плюс 20, Север минус 10. Второй способ: сдвиг 0, а районам пишем 120, 90 и 100. Прогнозы те же. Поэтому фразу «вес Центра» без договорённости сказать нельзя.

Жёлтая полоса, медленно. Сдвиг оставляем, один район выбрасываем. Вес оставшегося района~--- разница с выброшенным, не цена района вообще. В примере выброшен Юг, плюс 20 значит «Центр на 20 выше Юга».

Справа формула one-cold: \(c_{\text{ц}}+c_{\text{с}}+c_{\text{ю}}=2\mathbf 1\). В каждой строке две единицы, сумма столбцов равна двум столбцам сдвига. Та же ловушка. Столбец всё равно выбрасывают.

Словарь районов строят только по тем строкам, на которых учат. Район, который встретился только в проверке, своего столбца при обучении не получает.}

\slidehead{12}{Нормальные уравнения}
\screen{Слева жёлтая формула, три строки и матрица для прямой. Справа график: шаг за шагом ошибка падает к нулю.}
\begin{signs}
\item \(X\) — таблица признаков. Первый столбец — единицы.
\item \(\top\) сверху — транспонирование: строки и столбцы меняются местами.
\item \(X^{\top}X\) — признаки, помноженные друг на друга. Получается квадратная таблица.
\item \(w\) — столбец весов. Его ищем.
\item \(y\) — столбец настоящих ответов.
\item \(X^{\top}y\) — те же признаки, помноженные на ответы.
\item знак \(=\) в жёлтой формуле — веса стоят там, где сумма квадратов промахов меньше всего. Туда приходят шаги на картинке справа.
\item \(Q(w)\) — число, которое уменьшаем. Буква \(Q\) — просто имя этой суммы.
\item \(\|y-Xw\|^{2}\) — сумма квадратов промахов. \(y-Xw\) — промах каждой строки: ответ минус прогноз.
\item \(1/n\) — если сумму поделить на число объектов, место минимума не сдвинется. Вектор \(w\) останется тем же.
\item \(-2X^{\top}(y-Xw)\) — градиент, наклон этой суммы по весам. Минус двойка из производной квадрата.
\item «уравнение выше — это ноль градиента» — жёлтая формула получается, когда этот наклон равен нулю.
\item \(2X^{\top}X\) — гессиан, кривизна чаши. «Положительно полуопределён» значит: чаша открыта вверх, дно одно.
\item «один \(w\), только если столбцы независимы» — если столбцы копии, как районы вместе со сдвигом, весов много, прогноз один.
\item нижняя матрица, клетка \(n\) — число точек.
\item \(\sum x_i\) — сумма признака по всем точкам.
\item \(\sum x_i^{2}\) — сумма квадратов признака.
\item столбец \(a,\ b\) — сдвиг и наклон прямой. Это и есть искомые веса.
\item \(\sum y_i\) — сумма ответов.
\item \(\sum x_i y_i\) — сумма «признак умножить на ответ».
\end{signs}
\say{Сначала картинка справа. По горизонтали номер шага, по вертикали ошибка. На первом шаге ошибка большая. Каждый следующий шаг делает её меньше. В конце ошибка ноль. Так ищем веса.

Жёлтая формула слева, слева направо. \(X\)~--- таблица признаков, первый столбец единицы. Значок \(\top\)~--- строки стали столбцами. \(X^{\top}X\)~--- столбцы признаков, помноженные друг на друга. Дальше \(w\)~--- веса, их ищем. Справа от равно \(X^{\top}y\): та же таблица на столбец ответов \(y\). Равенство значит: веса стоят там, где сумма квадратов промахов меньше всего. Это место, куда пришли шаги. С обратной матрицы не начинай. Когда столбцы копии, обратной нет: это районы вместе со сдвигом.

Три строки под формулой. Читай буквы как написано.

Что минимизируем. \(Q(w)=\|y-Xw\|^{2}\). \(y-Xw\)~--- столбец промахов, настоящий ответ минус прогноз. Норма в квадрате~--- сумма квадратов этих промахов. Множитель \(1/n\) точку минимума не двигает.

Градиент. \(-2X^{\top}(y-Xw)\). Это наклон той же суммы по весам. В минимуме наклон ноль. Поставь скобку так, чтобы произведение стало нулём, сократи \(-2\)~--- получится жёлтая формула. На слайде так и написано: уравнение выше~--- это ноль градиента. Если вывод на доске не идёт, хватит этой фразы.

Гессиан. \(2X^{\top}X\), и написано: положительно полуопределён. Чаша открыта вверх, дно одно, минимум глобальный. Один набор \(w\) будет, только если столбцы не копии.

Нижняя матрица~--- та же жёлтая формула для прямой «старт \(a\) плюс наклон \(b\)». В квадрате слева: \(n\)~--- число точек, \(\sum x_i\)~--- сумма признака, \(\sum x_i^{2}\)~--- сумма квадратов признака. Столбец \(a,\ b\)~--- сдвиг и наклон. Справа \(\sum y_i\) и \(\sum x_i y_i\). На доске не решай. Числа 0,67 и 0,5 назовёшь на метриках, уже готовыми.

Если признака нет и в таблице один столбец единиц, формула схлопывается: константа равна среднему ответов на тех строках, где учили.}

\slidehead{13}{Когда формула молчит}
\screen{Три строки. Столбцы-копии. Почти копии. Компьютер не переворачивает матрицу.}
\begin{signs}
\item \(w\) — веса.
\item \(Xw\) — прогнозы всех строк. Разные \(w\) могут дать один и тот же \(Xw\): линия цен одна, весов много.
\item \(X^{\top}X\) — та квадратная таблица из жёлтой формулы. Её не переворачивают: переворот возводит ошибку округления в квадрат.
\item SVD и QR — имена разложения таблицы. Это способ решить систему, не переворачивая \(X^{\top}X\). Выводить не надо, достаточно назвать.
\end{signs}
\say{Прогноз может быть один, а весов много.

Копия столбца. На слайде: разные \(w\) дают один \(Xw\). \(Xw\)~--- прогнозы. Веса разные, линия цен одна. Это one-hot или one-cold вместе со сдвигом.

Почти копия. Площадь и число комнат ходят вместе. Маленький шум в ценах сильно качает оба веса, а линия цен почти стоит. «Рубли за метр» отдельно от «рублей за комнату» не произноси.

Как считают. На слайде написано SVD или QR: таблицу раскладывают, матрицу не переворачивают. Обращение \(X^{\top}X\) возводит ошибку округления в квадрат. Штраф, который чинит почти-копии, сегодня не выводим.}

\slidehead{14}{Квадрат и модуль}
\screen{Слева колокол: тогда минимум квадратов. Справа редкие большие промахи: тогда минимум модулей. Внизу: квадрат тянет прямую к дальней точке.}
\begin{signs}
\item \(y\) — столбец настоящих ответов.
\item \(Xw\) — прогноз прямой.
\item \(\varepsilon\) — промах, шум: ответ минус прогноз. Буква эпсилон, не «ошибка метрики».
\item \(\varepsilon_i\) — промах одной строки, номер \(i\).
\item «независимы» — промах одной квартиры не диктует промах следующей.
\item \(\mathcal{N}(0,\sigma^{2})\) — нормальный колокол. Первый аргумент \(0\) — среднее шума ноль. Второй, \(\sigma^{2}\) — разброс шума.
\item «максимум правдоподобия совпадает с минимумом квадратов» — при таком колоколе самое правдоподобное место весов и есть жёлтая формула.
\item справа, Лаплас — другое распределение шума. Его правдоподобие совпадает с минимумом модуля. Замкнутой жёлтой формулы у модуля нет.
\end{signs}
\say{Почему ошибку возводим в квадрат.

Слева карточка, формула \(y=Xw+\varepsilon\). \(y\)~--- ответы, \(Xw\)~--- прогноз прямой, \(\varepsilon\)~--- промах. Дальше: \(\varepsilon_i\) независимы и нормальные, \(\mathcal N(0,\sigma^{2})\), среднее ноль, разброс \(\sigma^{2}\). Это колокол. На слайде написано: максимум правдоподобия совпадает с минимумом квадратов. Поэтому у квадрата есть жёлтая формула.

Справа карточка «тяжёлые хвосты». На слайде: у распределения Лапласа правдоподобие совпадает с минимумом модуля. Замкнутой системы, как жёлтая формула, нет. Константу для модуля ставят в середину списка, в медиану.

Нижняя строка, медленно. Квадрат тянет прямую к дальней точке. Промах 10 входит в сумму как 100. Модуль смотрит, где лежит большинство: тот же промах входит как 10. Картинку откроем через несколько слайдов. Сплошная линия, обученная квадратом, уйдёт к жёлтой дальней точке. Пунктир, обученный модулем, останется у облака.}

\slidehead{15}{Четыре условия}
\screen{Четыре бытовые строки, когда весу можно верить. Жёлтая полоса: нормальный шум не нужен, причину формула не даёт.}
\say{Четыре строки. Это когда числу при признаке можно верить, а когда это просто то, что выплюнула формула. Название один раз: теорема Гаусса--Маркова. Слово BLUE не говори.

Среднее на прямой. Средняя цена и правда складывается из столбцов, которые мы положили: сдвиг, площадь, этаж. Если настоящая цена гнётся, а в таблице только прямая, веса описывают лучшую прямую по этим столбцам. Изгиб, которого в таблице нет, они не поймают.

Признак не из ошибки. Забыли метро. Близко к метро берут площадь побольше и платят больше. Часть денег за метро формула записывает в метры, вес при площади съезжает.

Промахи похожи. Примерно одного размера и не пачками. Если дорогие квартиры ошибаются на миллионы, а дешёвые на тысячи, прямая часто без сдвига в среднем, но уже не самая точная.

Столбцы разные. Площадь и число комнат ходят вместе. Вес «за метр» отдельно произнести нельзя.

Жёлтая полоса, медленно. Нормальный шум в этот список не входит. И формула не говорит, что площадь вызывает цену. Она говорит, как цена в среднем сдвигается вместе с площадью в этой таблице.

Если время поджимает, прочитай только жёлтую полосу и иди к метрикам.}

\slidehead{16}{Как читать один вес}
\screen{Плюс один к признаку сдвигает прогноз на его вес, если остальные столбцы не трогать. Слева что это значит. Справа чего это не значит.}
\begin{signs}
\item \(w_j\) — вес одного признака, номер \(j\).
\item «плюс один к признаку» — этот столбец увеличили на единицу.
\item «при прочих равных в модели» — остальные столбцы этой таблицы не трогаем. Сдвиг уже посчитан внутри выбранной записи.
\item после стандартизации единица — уже не метр, а один типичный разброс. Тогда \(w_j\) больше не «рубли за метр».
\end{signs}
\say{Фраза «плюс один к признаку сдвигает прогноз на его вес» значит вот что. Остальные столбцы этой таблицы зафиксированы. Сдвиг посчитан внутри той записи, которую выбрали.

Она не значит, что признак вызывает ответ. Не значит, что вес устойчив, если столбцы ходят вместе. И не значит, что после деления признака на его разброс вес всё ещё «рубли за метр». Тогда единица~--- один типичный разброс, не один метр.

Вес района~--- разница с выброшенным районом. В примере это было «Центр на 20 выше Юга».}

\slidehead{17}{Две работы}
\screen{Слева то, по чему ищут веса. Справа то, чем потом судят результат. Это разные числа.}
\say{Дальше много формул, и у них две разные работы.

Слева функция потерь. По ней на учебных строках ищут веса. Квадрат удобен: у него есть наклон и готовая формула.

Справа метрика. Ею судят уже найденные веса. Считать её надо на строках, которые в обучении не участвовали.

Ошибки 1 и 3. В квадрате это 1 и 9. В модуле это 1 и 3. Одно решение, два разных числа.

Учить квадратом, а в отчёт писать средний модуль~--- нормально, если так и сказано. Называть их одним числом нельзя.}

\slidehead{18}{MSE}
\screen{MSE~--- средняя квадратичная ошибка. Ошибки 1 и 3 входят как 1 и 9. На трёх точках около 0,06 у прямой и 0,22 у среднего. Справа чаша.}
\begin{signs}
\item \(\mathrm{MSE}\) — имя числа: средняя квадратичная ошибка.
\item \(m\) — сколько объектов в сумме. Не путать с \(n\) в матрице прямой: здесь буква из формулы на слайде.
\item \(1/m\) — делим сумму на число объектов, получаем среднее.
\item \(\sum_{j=1}^{m}\) — сложить слагаемые от первого объекта до последнего. \(j\) — номер объекта.
\item \(y_j\) — настоящий ответ объекта номер \(j\).
\item \(\hat y_j\) — прогноз для этого объекта.
\item \(y_j-\hat y_j\) — промах одного объекта.
\item квадрат снаружи скобки — промах возвели в квадрат. Ошибка 1 входит как 1, ошибка 3 как 9.
\item единицы — квадрат ответа. Если ответ в рублях, число в рублях в квадрате.
\end{signs}
\say{MSE читается так: средняя квадратичная ошибка. По-английски mean squared error. На слайде это написано под буквами.

Формула слева направо. \(\mathrm{MSE}\) равно. \(1/m\)~--- делим на число объектов \(m\). Знак суммы от \(j=1\) до \(m\)~--- складываем по всем объектам. В скобке \(y_j-\hat y_j\)~--- промах одного объекта, ответ минус прогноз. Квадрат снаружи скобки~--- этот промах в квадрате.

Ошибка 1 входит как 1. Ошибка 3 входит как 9. Далёкая точка получает большой голос. Если ответ в рублях, это число в рублях в квадрате. Неудобно, поэтому дальше будет корень.

Лучшая константа для этой ошибки~--- среднее учебных ответов. На наших трёх точках у прямой около 0,06, у фразы «всегда говори 1,67» около 0,22. Меньше лучше. Это всё ещё те же три строки.

Когда это говорят. Когда ищут веса: у квадрата есть формула, и большой промах должен весить сильнее маленького. В отчёт про цену квартиры это число не кладут. Оно в рублях в квадрате, клиент его не прочитает.}

\slidehead{19}{RMSE}
\screen{Корень из прошлой ошибки. Это не средняя ошибка. Среднюю ошибку в тех же единицах даёт следующий слайд.}
\begin{signs}
\item \(\mathrm{RMSE}\) — имя числа: корень из средней квадратичной ошибки.
\item знак корня — квадратный корень вокруг прошлого числа. Он возвращает рубли, если ответ был в рублях.
\item \(\mathrm{MSE}\) внутри корня — то число с прошлого слайда, уже с квадратами. Корень квадрат не отменяет: крупный промах всё равно вошёл сильнее, ещё до корня.
\end{signs}
\say{RMSE~--- корень из средней квадратичной ошибки. По-английски root mean squared error.

Формула короткая: \(\mathrm{RMSE}=\sqrt{\mathrm{MSE}}\). Корень из того, что только что посчитали. Квадрат уже сделал своё внутри MSE, до корня.

Число снова можно назвать в рублях. Крупная ошибка всё равно вошла через квадрат, ещё до корня. Корень это не отменяет.

Это не средняя ошибка. Среднюю ошибку в тех же единицах даёт следующий слайд.

Когда это говорят. Когда число хотят в рублях, но крупный промах всё равно должен весить больше. Так просят на соревновании, если в условии написан RMSE. Фраза вслух: типичный промах в рублях, но ошибка в миллион весит сильнее, чем десять ошибок по сто тысяч.}

\slidehead{20}{MAE}
\screen{Средняя абсолютная ошибка. Ошибки 1, 2 и 3 дают среднее 2. На трёх точках около 0,22 у прямой и 0,44 у среднего.}
\begin{signs}
\item \(\mathrm{MAE}\) — имя числа: средняя абсолютная ошибка.
\item \(1/m\) и \(\sum_{j=1}^{m}\) — те же, что в MSE: среднее по объектам.
\item вертикальные палочки \(|y_j-\hat y_j|\) — модуль промаха. Минус отбрасываем, квадрата нет. Ошибки 1, 2 и 3 входят как 1, 2 и 3.
\item единицы те же, что у ответа: рубли, если ответ в рублях.
\end{signs}
\say{MAE~--- средняя абсолютная ошибка. По-английски mean absolute error.

Формула как у MSE, только вместо квадрата модуль. \(1/m\)~--- снова среднее. Сумма по объектам. Внутри \(|y_j-\hat y_j|\)~--- модуль промаха. Единицы те же, что у ответа.

Ошибки 1, 2 и 3. Модули те же 1, 2 и 3. Среднее~--- 2. Квадрат из этой тройки сделал бы 1, 4 и 9.

Лучшая константа для модуля~--- середина списка, медиана. На трёх точках средний модуль у прямой около 0,22, у среднего около 0,44.

МНК~--- метод наименьших квадратов. Он искал веса по квадрату и не обещает победу по модулю. На этих трёх точках прямая выиграла и по модулю, но так бывает не всегда.

Когда это говорят. Когда спрашивают: на сколько рублей в среднем промахиваемся. Промах на сто тысяч и промах на миллион входят как сто тысяч и миллион, без квадрата. Это число для отчёта о цене.}

\slidehead{21}{Картинка хвоста}
\screen{Облако точек и одна дальняя жёлтая точка. Сплошная прямая ушла к ней, пунктир остался у большинства.}
\say{Жёлтая точка~--- хвост. Сплошная прямая обучена квадратом и ушла к ней. Пунктир обучен модулем и остался у большинства. Это картинка к двум прошлым слайдам. Линии нарисованы для глаза.}

\slidehead{22}{Доля разброса}
\screen{Доля разброса, которую забрала прямая. Внизу дроби~--- разброс самих ответов. Вверху~--- то, что прямая не угадала. Ноль значит «не лучше среднего».}
\begin{signs}
\item \(R^{2}\) — доля разброса. Не процент верных квартир. Буква \(R\) от слова «регрессия», двойка — «в квадрате», потому что внутри стоят квадраты.
\item \(1-\) перед дробью — из всего разброса вычитаем долю, которую прямая не забрала.
\item вверху дроби \(\sum(y_j-\hat y_j)^{2}\) — сумма квадратов промахов. Что прямая не угадала. Правый столбик.
\item внизу дроби \(\sum(y_j-\bar y)^{2}\) — сумма квадратов отклонений ответов от своего среднего. Разброс самих ответов. Левый столбик.
\item \(\bar y\) — среднее ответов. Чёрточка сверху, не шляпка. Шляпка была у прогноза.
\item \(y_j\) — ответ объекта номер \(j\).
\item \(\hat y_j\) — прогноз для него.
\item ноль у этой доли — прямая не лучше, чем просто сказать среднее. Не «ноль точности».
\end{signs}
\say{Эр-квадрат~--- доля разброса. На слайде так и написано: коэффициент детерминации. Это не процент верных ответов.

Формула слева направо. \(R^{2}\) равно единице минус дробь.

Внизу дроби \(\sum(y_j-\bar y)^{2}\). \(\bar y\)~--- среднее ответов. Каждый ответ минус это среднее, в квадрат, и всё сложить. Это разброс самих ответов. Левый столбик.

Вверху \(\sum(y_j-\hat y_j)^{2}\). Каждый промах в квадрат и сложить. Это то, что прямая не угадала. Правый столбик.

Дальше единица минус их отношение. Если правый столбик почти как левый, прямая почти ничего не забрала, доля около нуля. Если правый столбик низкий, прямая забрала почти весь разброс, доля около единицы. Подпись под столбиками уже посчитана: правый в четыре раза ниже, на слайде 0,75. Вслух дробь не пересчитывай.

Ноль значит «не лучше, чем просто сказать среднее». Это не «ноль точности». На новых данных доля бывает отрицательной: модель хуже, чем сказать среднее.

В программе берут среднее того ответа, который передали в функцию. Передали проверку~--- внизу стоит среднее проверки. В отчёте надо сказать, какое среднее использовали.

Когда это говорят. Когда хотят услышать, забрала ли прямая разброс цен по сравнению с фразой «всегда говори среднюю цену». Не когда хотят процент угаданных квартир.}

\slidehead{23}{Скорректированный}
\screen{Штраф за лишний столбец. Число строк и число столбцов вместе со сдвигом. Это не замена проверки на других строках.}
\begin{signs}
\item \(R^{2}_{\mathrm{adj}}\) — та же доля, но со штрафом за лишний столбец. Буквы adj — adjusted, «скорректированный».
\item \(R^{2}\) внутри — обычная доля с прошлого слайда.
\item \(1-R^{2}\) — доля, которую модель не забрала.
\item \(n\) — число строк, квартир.
\item \(p\) — число столбцов вместе со сдвигом. Сдвиг считается.
\item дробь \((n-1)/(n-p)\) — строк больше, чем столбцов, поэтому дробь больше единицы. Умножение на неё уменьшает долю.
\item итог \(1-\ldots\) — из единицы вычли уже увеличенный «незабранный» кусок. Поэтому скорректированная доля меньше обычной.
\end{signs}
\say{Обычная доля разброса на учебных строках почти обязана вырасти, если добавить столбец. Новый столбец не может сделать сумму промахов больше. Скорректированная доля этот бесплатный рост снимает.

Формула слева направо. \(R^{2}_{\mathrm{adj}}\) равно единице минус скобка. В скобке \((1-R^{2})\)~--- доля, которую модель не забрала. Её умножают на дробь \((n-1)/(n-p)\). \(n\)~--- число строк. \(p\)~--- число столбцов вместе со сдвигом. Строк больше, чем столбцов, поэтому дробь больше единицы, и итог меньше обычной доли.

На трёх точках строк 3, столбцов 2, обычная доля три четверти. Дробь равна 2. Одна четвёртая умножить на 2~--- одна вторая. Скорректированная доля падает с трёх четвертей до одной второй. Строк мало, столбцов два, штраф заметный.

Это штраф на тех же строках, где учили. Другие строки он не заменяет. Логарифм цены и саму цену так не сравнивают: разброс разный.

Когда это говорят. Добавили этаж или район, обычная доля на учебных строках выросла. Скорректированная показывает, не выросла ли она просто от нового столбца. Проверку на других строках она не заменяет.}

\slidehead{24}{MAPE}
\screen{Промах делят на сам ответ и усредняют. На трёх точках около 0,14 у прямой и 0,33 у среднего. При ответе 0 делить не на что.}
\begin{signs}
\item \(\mathrm{MAPE}\) — средняя абсолютная процентная ошибка. На слайде это доля, не «14 процентов».
\item \(1/m\) и сумма — снова среднее по объектам. \(m\) — сколько объектов, \(j\) — номер одного.
\item числитель \(|y_j-\hat y_j|\) — модуль промаха.
\item знаменатель \(|y_j|\) — модуль самого ответа. Делим промах на ответ, не на прогноз.
\item если \(y_j=0\), делить не на что. Такую строку формула не принимает.
\end{signs}
\say{MAPE~--- средняя абсолютная процентная ошибка. По-английски mean absolute percentage error.

Формула слева направо. \(1/m\) и сумма~--- снова среднее по объектам. В числителе \(|y_j-\hat y_j|\)~--- модуль промаха. В знаменателе \(|y_j|\)~--- модуль самого ответа. Промах делят на ответ.

На слайде это доля, не проценты. Около 0,14 у прямой и около 0,33 у среднего. Если кто-то скажет «четырнадцать процентов», поправь: на экране доля.

Маленький ответ весит больше большого, потому что на него делим. При ответе ноль делить не на что.

Когда это говорят. В одной таблице комнаты и дома. Промах на сто тысяч у комнаты за миллион и промах на два миллиона у дома за двадцать миллионов~--- оба около десяти процентов. Так сравнивают относительную ошибку. Не берут, если среди цен есть ноль.}

\slidehead{25}{sMAPE}
\screen{В знаменателе сумма модулей ответа и прогноза. Двойка сверху~--- договорённость.}
\begin{signs}
\item \(\mathrm{sMAPE}\) — та же процентная ошибка, но знаменатель другой. Буква s — symmetric, «симметричная». Слово про формулу, не про то, что оптимум несмещён.
\item \(1/m\) и сумма — среднее по объектам.
\item числитель \(2|y_j-\hat y_j|\) — двойка умножить на модуль промаха. Двойка — договорённость. Бывает запись без неё.
\item знаменатель \(|y_j|+|\hat y_j|\) — модуль ответа плюс модуль прогноза. Уже не делим на один ответ.
\item если и ответ, и прогноз нули, в знаменателе ноль. Такую долю договариваются отдельно.
\end{signs}
\say{sMAPE~--- симметричная средняя абсолютная процентная ошибка.

Формула слева направо. Снова \(1/m\) и сумма. В числителе \(2|y_j-\hat y_j|\): двойка на модуль промаха. Двойка~--- договорённость, встречается запись без неё. В отчёте надо назвать, какая именно. В знаменателе \(|y_j|+|\hat y_j|\): модуль ответа плюс модуль прогноза. Уже не делим на один ответ.

Слово «симметричная» описывает формулу. При том же модуле ошибки больший прогноз раздувает знаменатель и уменьшает штраф.

Если и ответ, и прогноз нулевые, в знаменателе ноль. Такую долю договариваются отдельно.

Когда это говорят. Когда ошибку хотят в процентах, а среди ответов есть очень маленькие. Продажи: вчера продали 2 штуки, прогноз 10. Обычная процентная ошибка делит промах только на 2 и становится огромной. Здесь в знаменателе 2 плюс 10, процент скромнее. Так сравнивают ряды, где рядом мелкие и крупные числа: продажи по дням, комнаты и дома в одной таблице. И помни нижнюю строку слайда: завышенный прогноз раздувает знаменатель, штраф становится меньше.}

\slidehead{26}{RMSLE}
\screen{Тот же корень, только от логарифмов. Ответ 100. Прогноз 50 дальше в логарифме, чем прогноз 150.}
\begin{signs}
\item \(\mathrm{RMSLE}\) — корень из средней квадратичной логарифмической ошибки. В конце буква L — logarithmic.
\item знак корня снаружи — как у RMSE.
\item \(1/m\) и сумма — среднее квадратов по объектам.
\item \(\log(1+y_j)\) — логарифм единицы плюс настоящий ответ.
\item \(\log(1+\hat y_j)\) — логарифм единицы плюс прогноз.
\item минус между логарифмами — сравниваем логарифмы, не рубли.
\item единица под логарифмом — чтобы нулевая цена не ушла в минус бесконечность.
\item квадрат снаружи скобки — разницу логарифмов возвели в квадрат, как в MSE.
\item отрицательный \(y\) в \(\log(1+y)\) не подставить.
\end{signs}
\say{RMSLE~--- корень из средней квадратичной логарифмической ошибки.

Формула слева направо. Снаружи корень, как у RMSE. Внутри \(1/m\) и сумма квадратов. В скобке под квадратом \(\log(1+y_j)-\log(1+\hat y_j)\): логарифм единицы плюс ответ минус логарифм единицы плюс прогноз. Единица под логарифмом нужна, чтобы нулевая цена не ушла в минус бесконечность. Сравниваем не рубли, а логарифмы.

Пример. Ответ 100. Прогноз 50 и прогноз 150 ошибаются на одни и те же 50 рублей. В логарифме занижение даёт около 0,68, завышение около 0,40. Вслух достаточно: занизили~--- штраф больше. Метрика смотрит, во сколько раз промахнулись.

Отрицательный ответ в логарифм не подставить. Если бизнесу нужны рубли, это число ему рублёвую ошибку не пересказывает.

Когда это говорят. Когда важно, во сколько раз промахнулись, и занижение дороже завышения. Цены от миллиона до десятков миллионов. Занизили вдвое~--- штраф больше, чем завысили вдвое. Не берут, если в отчёте нужны именно рубли.}

\slidehead{27}{Pinball}
\screen{Штраф за недобор и перебор разный. При уровне одна вторая это половина модуля. Справа график скошен.}
\begin{signs}
\item \(\rho_{\tau}(r)\) — штраф одного промаха. Буква ро — просто имя функции. Не аббревиатура.
\item \(\tau\) — какую долю списка хотим накрыть. Число от 0 до 1. Тау 0,8 — отметка, ниже которой 80 цен из 100.
\item \(r\) — остаток: ответ минус прогноз. Плюс значит прогноз не добрал. Минус значит прогноз перебрал.
\item верхняя строка фигурной скобки, \(r\ge 0\) — случай недобора. Штраф равен \(\tau r\).
\item нижняя строка, \(r<0\) — случай перебора. Штраф равен \((\tau-1)r\). Скобка \(\tau-1\) отрицательная, а \(r\) тоже отрицательный, поэтому штраф положительный и меньше, чем за недобор, если \(\tau\) больше половины.
\item при \(\tau=1/2\) обе ветки дают половину модуля. Лучшая константа тогда — медиана.
\end{signs}
\say{Pinball~--- не аббревиатура. Это штраф для квантиля, то есть для выбранной доли списка.

Формула слева направо. \(\rho_{\tau}(r)\)~--- штраф. Буква \(\tau\)~--- какую долю списка хотим накрыть. \(r\)~--- остаток, ответ минус прогноз. Справа фигурная скобка, две строки.

Верхняя строка: если \(r\ge 0\), штраф равен \(\tau r\). Положительный остаток значит: прогноз не добрал.

Нижняя строка: если \(r<0\), штраф равен \((\tau-1)r\). Прогноз перебрал.

Если не добрали, штраф больше, когда \(\tau\) больше половины. Если перебрали, штраф меньше. При \(\tau=1/2\) это половина модуля, и лучшая константа~--- медиана. На слайде внизу: оптимум по константе~--- квантиль уровня \(\tau\), не среднее.

Когда это говорят. Нужна не средняя цена, а осторожная отметка. Тау 0,8 значит: хотим цену, ниже которой окажутся 80 цен из 100. Настоящая цена окажется выше прогноза только в 20 случаях из 100. Недобор при этом штрафуется сильнее.}

\slidehead{28}{Хьюбер}
\screen{Рядом с нулём штраф как у квадрата, дальше растёт как у модуля. Это способ учить, не число для отчёта.}
\begin{signs}
\item \(H_{\delta}(r)\) — штраф одного промаха. \(H\) — фамилия Хьюбера, не аббревиатура.
\item \(r\) — промах: ответ минус прогноз.
\item \(\delta\) — порог. Маленький порог ближе к модулю, большой — к квадрату. Его выбирают на отложенных строках.
\item верхняя строка, \(|r|\le\delta\) — промах не дальше порога. Штраф \(\tfrac{1}{2}r^{2}\): половина квадрата. Рядом с нулём функция гладкая.
\item нижняя строка, \(|r|>\delta\) — промах дальше порога. Штраф \(\delta(|r|-\delta/2)\): растёт как модуль, не как квадрат. Хвост не получает квадратичного голоса.
\item вертикальные палочки — модуль промаха.
\end{signs}
\say{Хьюбер~--- фамилия, не аббревиатура.

Формула слева направо. \(H_{\delta}(r)\)~--- штраф промаха \(r\). Буква \(\delta\)~--- порог. Справа две строки.

Верхняя: если модуль промаха не больше \(\delta\), штраф равен \(\tfrac12 r^{2}\). Рядом с нулём это квадрат, функция гладкая.

Нижняя: если модуль больше \(\delta\), штраф равен \(\delta\) умножить на скобку \(|r|-\delta/2\). Дальше порога рост линейный, как у модуля. Хвост не получает квадратичного голоса.

Маленький порог ближе к модулю, большой~--- к квадрату. Порог выбирают на отложенных строках. На слайде внизу: это способ учить, не число в отчёт. В отчёт пишут средний модуль или корень из квадрата.

Когда это говорят. Учим прямую, а в таблице несколько очень дорогих квартир. Квадрат утащит линию к ним. Модуль негладкий, готовой формулы нет. Хьюбер в середине как квадрат, на хвосте как модуль. Клиенту это число не называют.}

\slidehead{29}{Учебные и проверочные строки}
\screen{Три кучи. Учебные: веса и средние. Отложенные: какой набор столбцов оставить. Совсем отложенные: один отчёт. Жёлтая полоса: среднее и словарь районов только с учебных строк.}
\begin{signs}
\item \(\mu\) на строке train — среднее признака, посчитанное только на учебных строках. Им потом сдвигают и проверку.
\item \(\sigma\) там же — разброс признака, тоже только с учебных строк. Проверку на это среднее и разброс не пересчитывают.
\end{signs}
\say{Числа квадрата, модуля и доли разброса, которые только что называли, посчитаны на тех же строках, где искали веса. Это ещё арифметика. Качество считают на других строках.

На учебных строках ищут веса. Туда же среднее константы, словарь районов, среднее и разброс для масштаба. Мю и сигма на слайде~--- это среднее и разброс.

Отложенные строки выбирают, какой набор столбцов оставить, какую метрику писать в отчёт и какой порог у Хьюбера.

Самые отложенные строки~--- один отчёт уже выбранного решения. Поправили модель после взгляда на них~--- отчёт уже нечестный.

Жёлтая полоса. Среднее, масштаб и словарь районов считают только на учебных строках.}

\slidehead{30}{Четыре разреза}
\screen{Четыре картинки. Случайный: цвета вперемешку. Корзины цены: две кучи, в каждой есть дорогие с жёлтой рамкой. Время: слева прошлое, справа будущее. Группа: дом целиком одним цветом.}
\say{Как разрезать таблицу, чтобы проверка была похожа на будущий прогноз. Светлый квадрат ушёл в учебные строки, по нему ищут веса. Тёмный ушёл в проверку, по нему только проверяют.

Случайный. Квадраты перемешаны. Так режут, когда строки ничем не связаны: разные квартиры, нет общего дома и нет порядка по дням.

Корзины цены. Здесь две кучи. В учебной куче светлые квадраты: часть обычные, часть с жёлтой рамкой. В проверочной куче тёмные: тоже часть обычные, часть с жёлтой рамкой. Жёлтая рамка~--- дорогая квартира. Она не отдельная куча. Дорогие лежат и там, и там. Корзины режут заранее: дешёвые, средние, дорогие. Саму цену ножом разреза не режут: цена~--- это ответ.

Время. Слева прошлое, все квадраты светлые. Справа будущее, все тёмные. Не перемешивают. Завтрашняя цена не имеет права учить вес для вчера.

Группа. Квадраты одного дома стоят одним цветом: либо весь дом в учебе, либо весь дом в проверке. Если разрезать дом по сделкам, прямая запомнит дом.

Честный разрез тот, который повторяет будущий прогноз. Цена на завтра требует времени. Прогноз по новому дому требует группы.}

\slidehead{31}{Где разрез ломается}
\screen{Четыре поломки. Не перемешали. Время перемешали. Дом разрезали по строкам. Словарь районов собрали до разреза.}
\say{Четыре способа сломать разрез. Во всех уравнение видит строки, которых в бою не будет.

Не перемешали. Таблица отсортирована по цене, кусок отрезали как есть. Проверка оказывается другим режимом цен.

Время перемешали. Завтрашняя цена стоит в учебе и учит вес для вчера.

Дом разрезали по строкам. Один дом и в учебе, и в проверке. Прямая запоминает дом.

Словарь районов собрали до разреза. Столбцы собраны по всей таблице, включая районы, которые есть только в проверке. Район только из проверки своего веса не получает.}

\slidehead{32}{Три утечки}
\screen{В таблицу попало то, чего не будет в момент прогноза. Прокси ответа. Статистика до разреза. Один дом по обе стороны.}
\say{В таблицу попало то, чего не будет в момент прогноза.

Прокси ответа. Цена за метр, умноженная на площадь, восстанавливает цену. Веса будут красивыми и пустыми: модель увидела ответ в признаке.

Статистика до разреза. Среднее пропуска, среднее и разброс для масштаба, словарь районов~--- только с учебных строк. Если спросят про числа: в учебе 0, 0 и 10, в проверке 10. Среднее учебы~--- около 3,3. Среднее всех четырёх~--- 5. Второе уже видело проверку.

Одна сущность дважды. Сделки одного дома по обе стороны. Модель запоминает дом.}

\slidehead{33}{Проверка}
\screen{Вопрос: какие утверждения верны. А~--- шаги, ошибка меньше. Б~--- минимум суммы квадратов. В~--- коды 1, 2, 3. Ответ по клику: верно А и Б.}
\say{Минута в чат. Какие утверждения верны. Пока не открывай ответ.

А. Веса ищем шагами, на каждом шаге ошибка меньше. Верно. Это график рядом с жёлтой формулой.

Б. Веса стоят там, где сумма квадратов промахов меньше всего. Тоже верно. Шаги ведут именно туда. Жёлтая формула записывает это место.

В. Район кодируют числами 1, 2, 3. Неверно. Модель прочитает Юг как втрое больший Центр. Район превращают в столбцы и один уровень выбрасывают.

Открывай. Верно А и Б.}

\slidehead{34}{Итог}
\screen{Шесть коротких строк. Прочитай список как есть.}
\say{На этом всё. Читай список как есть.

Прогноз~--- это сдвиг плюс веса признаков.

Веса ищем шагами, ошибка падает, формула записывает минимум.

Район превращаем в столбцы и один уровень выбрасываем.

Смотрели квадрат ошибки, модуль и долю разброса.

Учат на учебных строках. Проверяют на других.

Первого ноября открываете задачу и смотрите, ответ там число или класс. Если класс, модель строим со следующей пары.}

\end{document}
